% !TeX root = main.tex
\section{Introduction}

The concept of a \emph{twisted sheaf} goes
back to Giraud 
and his formal treatment of the Brauer group in terms of \emph{gerbes}~\cite[Section~V.4]{giraud1971}.
During the last two decades,
twisted sheaves and their derived categories have gained a renewed interest,
starting with the thesis of C{\u{a}}ld{\u{a}}raru \cite{caldararu2000},
who studied moduli problems for semi-stable sheaves on varieties,
and the work by de~Jong~\cite{dejong2004}
and Lieblich~\cite{lieblich2004, lieblich2015}
on the period--index problem for the Brauer group.
For a survey on later developments,
we refer to Lieblich~\cite[Section~4]{lieblich2017}.

Although twisted sheaves have a rather elementary description in terms of 2-cocycles
(cf.~Remark~\ref{rem-twisted-2-cocycle}),
they are most naturally viewed as sheaves on gerbes banded by diagonalizable groups
(see~Definition~\ref{def:twisted}).
Our first result is the following theorem on the decomposition of the derived category associated to such a gerbe,
generalizing a result by Lieblich~\cite[Section~2.2.4]{lieblich2004}.

\begin{mtheorem}[{see Theorem~\ref{t:chi-splitting-D-qc-X}}]
  \label{t:chi-splitting-D-qc-X-intro}
  Let $S$ be an algebraic stack and 
  $\pi\colon X \to S$ a gerbe banded by a diagonalizable group
  $\Delta$ 
  with character group $A$.
  Then taking the coproduct defines an equivalence 
  \begin{equation}
    \label{eq-equivalence-derived-intro}
    \prod_{\chi \in A} \D_{\qc, \chi}(X) \sira \D_\qc(X),
    \qquad
    (\mathcal{F}_\chi)_{\chi \in A} 
    \mapsto 
    \bigoplus_{\chi \in A} \mathcal{F}_\chi,
  \end{equation}
  of triangulated categories,
  where $\D_{\qc, \chi}(X)$ denotes the full subcategory of
  $\D_\qc(X)$ of objects with $\chi$-homogeneous quasi-coherent
  cohomology (see Definition~\ref{def-homogeneous-sheaf}).
\end{mtheorem}

In the underived setting, 
a splitting similar as in the theorem above (see Theorem~\ref{t-qc-gerbe-split})
is a consequence of the well-known fact that quasi-coherent
representations of $\Delta$ split into subrepresentations
corresponding to the characters of the group~$\Delta$.

The result by Lieblich is the decomposition in Theorem~\ref{t:chi-splitting-D-qc-X-intro}
in the special case when $S$ is a quasi-compact, separated scheme.
He proves and uses the fact that in this setting the obvious functor
\begin{equation}
\label{eq-dqcoh-vs-dqcoh}
\D(\Qcoh(X)) \to \D_\qc(X)
\end{equation}
is an equivalence of categories~\cite[Proposition~2.2.4.6]{lieblich2004}.
Note that, if the base $S$ is a scheme, the functor \eqref{eq-dqcoh-vs-dqcoh} only fails to be an equivalence
in truly pathological situations.
However, as shown by Hall--Neeman--Rydh \cite[Theorem~1.3]{hnr2018},
the functor \eqref{eq-dqcoh-vs-dqcoh} fails to be an equivalence for large classes of algebraic stacks,
including the basic case when $X$ is the classifying stack $\BB \GL_n$ of the general linear group
(see Remark~\ref{rem-trivial-case}).
In particular, this kind of argument cannot be used to obtain the splitting of the derived category
of one of the most fundamental examples of gerbes banded by $\GGm$ ---
namely the \emph{gerbe of trivializations}
$\BB \GL_n \to \BB \PGL_n$ for the universal Brauer--Severi scheme
(see Remark~\ref{rem-universal-gerbe-of-trivializations}).

In order to prove Theorem~\ref{t:chi-splitting-D-qc-X-intro},
we study the more general problem whether suitable torsion pairs in abelian categories induce
semi-orthogonal decompositions on the level of derived categories.
We achieve the following theorem of independent interest.

\begin{mtheorem}[{see Theorem~\ref{t:sod-DbBA} and Theorem~\ref{t:sod-DBA}}]
  \label{t:sod-DBA-intro}
  Let $\mathcal{B}$ be a weak Serre subcategory of an abelian
  category $\mathcal{A}$.
  Let 
  $\mathcal{T}$ and
  $\mathcal{F}$ be abelian subcategories of $\mathcal{B}$
  forming a torsion pair $(\mathcal{T}, \mathcal{F})$ in
  $\mathcal{B}$.
  Assume that 
  \begin{equation*}
    \Ext^n_{\mathcal{A}}(T, F)=0
  \end{equation*}
  for all objects $T \in \mathcal{T}$, $F \in \mathcal{F}$ and
  all integers $n$. 
  Then 
  \begin{equation*}
    \D_{\mathcal{B}}^\bd(\mathcal{A})=
    \langle \D_{\mathcal{F}}^\bd(\mathcal{A}), 
    \D_{\mathcal{T}}^\bd(\mathcal{A})\rangle
  \end{equation*}
  is a semi-orthogonal decomposition.
  Furthermore, if the inclusion $\D_{\mathcal{B}}(\mathcal{A})
  \subset \D(\mathcal{A})$ 
  satisfies some technical
  conditions on homotopy limits and colimits
  formulated precisely in the statement of Theorem~\ref{t:sod-DBA},
  then
  \begin{equation*}
    \D_{\mathcal{B}}(\mathcal{A})=
    \langle \D_{\mathcal{F}}(\mathcal{A}), 
    \D_{\mathcal{T}}(\mathcal{A})\rangle
  \end{equation*}
  is a semi-orthogonal decomposition.
\end{mtheorem}

The theory of gerbes,
hinging on the theory of stacks,
used to have a reputation of being inaccessible,
and Giraud's book on non-abelian cohomology \cite{giraud1971} introducing them is notorious for being a hard read,
even among experts.
However,
in recent years the theory of algebraic stacks has become a main stream part of algebraic geometry,
much owing to the excellent text books by Laumon--Moret-Bailly~\cite{lmb2000} and Olsson~\cite{olsson2016},
and to the Stacks Project~\cite{stacks-project}.
Moreover,
as explained in the thesis by Lieblich \cite{lieblich2004},
only a small part of Giraud's theory is actually needed to
develop a satisfactory theory of twisted sheaves.
We illustrate the effectiveness of the language of gerbes by giving a short,
simple and conceptually appealing proof of the following theorem,
which generalizes a result by Bernardara \cite[Theorem~4.1]{bernardara-BS-schemes}.

\begin{mtheorem}[{see Theorem~\ref{t:sod-brauer-severi} and
    Corollary~\ref{c:sod-brauer-severi}}]
  \label{t:sod-brauer-severi-intro}
  Let $S$ be an algebraic stack and $\pi\colon P \to S$ a
  Brauer--Severi 
  scheme of relative dimension $n$ over $S$.
  Let $\beta \in \HH^2(S_\fppf, \GGm)$ denote the Brauer class of
  $\pi$. 
  Then, for any $a \in \bZ$, the category $\D_\qc(P)$ admits a
  semi-orthogonal 
  decomposition
  \begin{equation}
    \label{eq:SOD-Dpf}
    \D_\qc(P) = \langle \mathcal{D}_a, \ldots, \mathcal{D}_{a+n}
    \rangle 
  \end{equation}
  into right admissible subcategories $\mathcal{D}_i$,
  where the category $\mathcal{D}_i$ is equivalent to 
  the category $\D_{\qc,\beta^i}(S)$
  of $\beta^i$-twisted complexes on $S$ (see
  Definition~\ref{def:derived-twisted} and Remark~\ref{r:image-twisted}). 
  
  Similarly, we have corresponding decompositions
  for the category of perfect complexes, the category of locally bounded pseudo-coherent complexes
  and the singularity category.  
  In particular, we also have such a decomposition for $\D^b_\co(P)$ when
  $S$ is Noetherian.
\end{mtheorem}

Bernardara's result is the semi-orthogonal decomposition 
\eqref{eq:SOD-Dpf} in the case when $S$ is a Noetherian, separated
scheme with the property that any pair of points in $S$ is
contained in an open affine subscheme (cf.\ Remark~\ref{r:bernardara}).
Apart from our Theorem~\ref{t:chi-splitting-D-qc-X-intro},
the main ingredient in our proof
of Theorem~\ref{t:sod-brauer-severi-intro} is a classical
result by Orlov regarding a semi-orthogonal decomposition of the derived
category of a projectivized vector bundle \cite[Theorem~2.6]{orlov1992},
which we previously have generalized to algebraic stacks using the technique of \emph{conservative descent}
\cite[Theorem~6.7, Corollary~6.8]{BS-conservative}.

While finishing the manuscript, we were informed by Brown and~Moulinos of a result \cite[Theorem 3.1]{bm2019}
which is similar to our Theorem~\ref{t:sod-brauer-severi-intro}.

\subsection*{Outline}
In Section~\ref{sec:torsion-pairs-decomp},
we work mostly in the general setting of derived categories of abelian categories.
This section contains the proof of Theorem~\ref{t:sod-DBA-intro}.
In Section~\ref{sec:gerbes-and-twisted} and Section~\ref{sec:quasi-coher-sheav},
we recall and collect some basic facts on gerbes,
bandings and twisted sheaves which are scattered in the literature.
We expect most of these results to be well known to experts.
In Section~\ref{sec:decomp-deriv-categ}, we turn our attention to derived categories of gerbes and twisted sheaves.
We generalize some of the fundamental results by C{\u{a}}ld{\u{a}}raru and Lieblich \cite{caldararu2000, lieblich2004}
to the case when we work over an algebraic stack.
Most notably we prove Theorem~\ref{t:chi-splitting-D-qc-X-intro}.
Finally, in the last section,
we give a brief summary on Giraud's treatment of the Brauer group \cite[Section~V.4]{giraud1971},
and use the theory to prove Theorem~\ref{t:sod-brauer-severi-intro}.

\subsection*{Acknowledgments}
The first named author was partially supported by the Danish National
Research Foundation through the Niels Bohr Professorship of Lars Hesselholt,
by the Max Planck Institute for Mathematics in Bonn,
and by the DFG through SFB/TR 45.
The second named author was partially supported by the DFG through a postdoctoral fellowship and through SFB/TR 45.

\subsection*{Notation and conventions}
\label{sec:conventions}

If $\mathcal{A}$ is an abelian category, we write
$\D(\mathcal{A})$ for its unbounded derived category. 
We write $\D^{\leq n}(\mathcal{A})$ and 
$\D^{\geq n}(\mathcal{A})$ for its full subcategories of objects
whose cohomology is concentrated in the indicated degrees, where
$n \in \bZ$.
The (intelligent) truncation functors on $\D(\mathcal{A})$ are
denoted $\tau_{\leq n}$ and $\tau_{\geq n}$. 
If $E$ and $F$ are objects of $\D(\mathcal{A})$, we 
sometimes abbreviate
$\Ext^n_\mathcal{A}(E,F):=\Hom_{\D(\mathcal{A})}(E, \Sigma^nF)$.

By a \emph{diagonalizable group} we mean a diagonalizable group scheme
over $\Spec \bZ$ that is isomorphic to the spectrum of the group
ring of a finitely generated abelian group.

We use the definitions of algebraic space and algebraic stack given in the stacks project
\cite[\sptag{025Y}, \sptag{026O}]{stacks-project}.
In particular, we do not impose any separatedness conditions. 
The algebraic stacks form a 2-category where each 2-morphism is invertible.
We will follow the common practice to often suppress 2-categorical details.
For instance, we will usually simply write morphism of stacks instead of 1-morphism
of stacks, subcategory of the category of stacks instead of 2-subcategory etc.

For the theory of quasi-coherent modules and derived categories in the context of algebraic stacks,
we basically follow the approach of Laumon--Moret-Bailly~\cite[Sections~12 and~13]{lmb2000},
which we now briefly recapitulate.
Let $X$ be an algebraic stack.
We denote the topos of sheaves on the big \emph{fppf} site over $X$ by $X_\fppf$
and the topos of sheaves on the \emph{lisse--étale} site over $X$ by $X_\liset$.
For $\tau \in \{\fppf, \liset\}$,
we let $\Mod(X_\tau, \mathcal{O}_X)$ denote category of $\mathcal{O}_X$-modules in $X_\tau$,
and $\Qcoh(X_\tau, \mathcal{O}_X)$ the full subcategory of quasi-coherent modules
(see~\cite[\sptag{03DL}]{stacks-project}).
\begin{remark}
\label{rem-coherator}
The categories $\Qcoh(X_\tau, \mathcal{O}_X)$ and $\Mod(X_\tau, \mathcal{O}_X)$ are Grothendieck abelian categories
\cite[\sptag{0781}]{stacks-project}.
In particular,
since the inclusion functor $\Qcoh(X_\tau, \mathcal{O}_X) \hookrightarrow \Mod(X_\tau, \mathcal{O}_X)$ is colimit preserving,
it admits a right adjoint.
This right adjoint is called the \emph{quasi-coherator}.
\end{remark}

The construction $(-)_\tau$ functorially associates an adjoint pair 
\begin{equation}
\label{eq-functorial-topos}
f^*\colon Y_\tau \rightleftarrows X_\tau \colon f_*
\end{equation}
of functors to any morphism $f\colon X \to Y$ of algebraic stacks.
Moreover, the functor $f^*$ preserves finite products.
In particular, the adjunction \eqref{eq-functorial-topos} induces an adjunction
\begin{equation}
\label{eq-functorial-module}
f^*\colon \Mod(Y_\tau, \mathcal{O}_Y) \rightleftarrows \Mod(X_\tau, \mathcal{O}_X) \colon f_*.
\end{equation}
Moreover, the pull-back functor $f^*$ preserves quasi-coherence.
The functor
\begin{equation}
\label{eq-qc-push-forward}
f_*\colon \Qcoh(X_\tau, \mathcal{O}_X) \to \Qcoh(Y_\tau, \mathcal{O}_Y)
\end{equation}
is defined as the right adjoint of the restriction of $f^*$ to
quasi-coherent sheaves.
The existence of such a right adjoint is guaranteed by the quasi-coherator
(see~Remark~\ref{rem-coherator}).
\begin{remark}
\label{rem-qc-pf}
If $\tau = \liset$, then the functor \eqref{eq-qc-push-forward} is the restriction of the push-forward for sheaves of modules
provided that $f$ is quasi-compact and quasi-separated,
but not in general.
If $\tau = \fppf$, then \eqref{eq-qc-push-forward} seldom is the restriction of the push-forward for sheaves of modules.
\end{remark}
The category $\Mod(X_\tau, \mathcal{O}_X)$ has a closed symmetric monoidal
structure with operations given by the usual tensor product  and the sheaf hom functor.
The tensor operation preserves quasi-coherence and is preserved by pull-backs.
In particular, we get an induced symmetric monoidal structure on $\Qcoh(X_\tau, \mathcal{O}_X)$.

\begin{remark}
\label{rem-qc-hom}
The symmetric monoidal structure on $\Qcoh(X_\tau, \mathcal{O}_X)$ induced by the tensor product
is also closed.
This follows from the existence of the quasi-coherator (see~Remark~\ref{rem-coherator}).
For $\tau = \liset$,
the internal hom functor coincides with the sheaf hom functor provided that the first argument is of finite presentation, but not in general.
\end{remark}

There is an obvious restriction functor $X_\fppf \to X_\liset$,
which is compatible with push-forward.
This induces a monoidal equivalence
\begin{equation}
\label{eq-change-site}
\Qcoh(X_\fppf, \mathcal{O}_X) \sira \Qcoh(X_\liset, \mathcal{O}_X),
\end{equation}
functorial in $X$ (see~\cite[\sptag{07B1}]{stacks-project}).

We denote the derived category $\D(\Mod(X_\liset, \mathcal{O}_X))$ simply by $\D(X)$.
The category $\D_\qc(X)$ is defined as the full subcategory of $\D(X)$ consisting of
complexes with quasi-coherent cohomology sheaves.
The derived tensor product on $\D(X)$ preserves objects in $\D_\qc(X)$,
giving the latter category a symmetric monoidal structure, which is closed.
\begin{remark}
\label{rem-dqc-hom}
The internal hom functor on $\D_\qc(X)$ is the restriction of the derived sheaf hom functor provided
that the first argument is perfect, but not in general
(see~\cite[Section~1.2]{hr2017}). 
\end{remark}
Given an arbitrary morphism $f\colon X \to Y$ of algebraic stacks,
we have an adjoint pair of functors
\begin{equation}
\label{eq-pb-pf}
f^*\colon \D_\qc(Y) \to \D_\qc(X), \qquad f_*\colon \D_\qc(Y) \to \D_\qc(X),
\end{equation}
and the formation of such pairs is functorial.
\begin{remark}
\label{rem-derived-pb}
The functor $f^*$ in \eqref{eq-pb-pf} coincides with the derived pull-back between the ambient categories provided that
$f$ is smooth.
For general $f$, the construction of $f^*$ is somewhat technical owing to the fact that the adjunction
\eqref{eq-functorial-topos}, for $\tau = \liset$,
does not form a morphism of topoi,
as was first noted by Behrend~\cite[Warning~5.3.12]{behrend-derived-l-adic-stacks}.
The construction is worked out by Olsson~\cite{olsson2007} in the bounded case and by Laszlo--Olsson~\cite[Example~2.2.5]{lo2008} in the unbounded case.
We refer to \cite[Section~1]{hr2017} for a concise summary on this.
\end{remark}
\begin{remark}
\label{rem-derived-pf}
The push-forward $f_*$ in \eqref{eq-pb-pf} coincides with the derived push-forward provided that
$f$ is \emph{concentrated} \cite[Definition~2.4, Theorem~2.6(2)]{hr2017},
but not in general.
\end{remark}

\begin{remark}
It is also possible to construct $\D_\qc(X)$ using the topos $X_\fppf$.
This approach is taken in the stacks project \cite[\sptag{07B6}]{stacks-project}.
For $\tau = \fppf$ the adjunction \eqref{eq-functorial-topos} does give a morphism of topoi,
making the construction of the pull-back in \eqref{eq-pb-pf} easier.
On the other hand the inclusion $\Qcoh(X_\fppf, \mathcal{O}_X) \subset \Mod(X_\fppf, \mathcal{O}_X)$
is not exact \cite[\sptag{06WU}]{stacks-project},
making the actual construction of $\D_\qc(X)$ quite technical.
We will not use this point of view in the present article.
\end{remark}

\begin{notation}
\label{not-sheaf-op}
We conclude this section by summarizing our notational conventions regarding sheaves on stacks.
In the discussion below, we let $\tau \in \{\fppf, \liset\}$.
\begin{enumerate}
\item
We simply write $\Qcoh(X)$ for any of the categories $\Qcoh(X_\tau, \mathcal{O}_X)$
when no distinction is necessary.
This is motivated by the functorial equivalence~\eqref{eq-change-site}.
\item
The monoidal operations on the categories
$$
\Qcoh(X_\tau, \mathcal{O}_X), \qquad
\Mod(X_\tau, \mathcal{O}_X), \qquad
\D(X), \qquad
\D_\qc(X)
$$
are always denoted $-\otimes-$ and $\sheafHom(-,-)$, respectively.
In particular, we omit any derived decorations.
The precise category we are working on will always be inferable from context.
Note that in the cases $\Qcoh(X_\tau, \mathcal{O}_X)$ and~$\D_\qc(X)$
the functor $\sheafHom$ does not always coincide with $\sheafHom$ on
the ambient category (cf.~Remark~\ref{rem-qc-hom} and~Remark~\ref{rem-dqc-hom}). 
\item
The pull-backs and push-forwards along a morphism $f\colon X \to Y$ are always denoted by $f^*$ and~$f_*$,
respectively.
Which of the categories
$$
(-)_\tau,\qquad
\Mod((-)_\tau, \mathcal{O}_{(-)}), \qquad
\Qcoh((-)_\tau, \mathcal{O}_{(-)}),
$$
$$
\D(-), \qquad
\D_\qc(-),
$$
we are working with will always be inferable from context.
In particular, 
we do not use any derived decorations or any particular decorations for
push-forward of quasi-coherent sheaves (cf.~Remark~\ref{rem-qc-pf} and~Remark~\ref{rem-derived-pf}).
\end{enumerate}
\end{notation}

%%% Local Variables: 
%%% mode: latex
%%% TeX-master: "main"
%%% End: 
